\subsection{庞加莱群的容许拓扑}

设 B 是拓扑空间，\(a\) 是 B 的一点。称 \(\pi_1(B, a)\) 的子群 H 为容许的，如果 B 的每一点 \(b\) 都有邻域 V，使得 \(\gamma i_*(\delta)\gamma^{-1} \in \mathrm{H}\) 对每个 \(\gamma \in \varpi_{a,b}(\mathrm{B})\) 和每个 \(\delta \in \pi_1(V, b)\) 都成立，其中 \(i: V \to B\) 是典范嵌入。如果 H 是 \(\pi_1(B, a)\) 的正规子群，则要使 H 容许，只需对一个道路类 \(\gamma \in \varpi_{a,b}(\mathrm{B})\) 验证此条件即可。

命题 4. — \(\pi_1(B, a)\) 上存在唯一一个与其群结构相容的拓扑，使 \(\pi_1(B, a)\) 的容许正规子群构成单位元的邻域基。在此拓扑下，开子群恰好是容许子群。

由容许子群的定义得到以下事实：

\begin{enumerate}
\def\labelenumi{\alph{enumi})}
\item
  群 \(\pi_1(B, a)\) 是容许的；
\item
  包含某个容许子群的子群是容许的；
\item
  有限个容许子群的交是容许的。
\item
对于 \(\pi_{1}(B,a)\) 的每个容许子群 H，各子群 \(\gamma H\gamma^{-1}\) 的交（其中 \(\gamma\) 遍历 \(\pi_{1}(B,a)\)）是容许的。
\end{enumerate}

特别地，容许正规子群的集合是一个滤基，由 \(\pi_1(B,a)\) 中满足公理 (GV \(_{\text{III}}'\) ) 的子群组成；根据 TG，III，第 5 页，例，这就得到命题的第一部分。第二部分随即显然（参见 TG，III，第 7 页，命题 4 的推论）。

命题 4 所刻画的群 \(\pi_{1}(B,a)\) 上的拓扑称为容许拓扑。

注记。--- 1) 若 B \(_{0}\) 表示 B 中 \(a\) 的道路连通分支，则在这些群赋予容许拓扑时，典范同构 \(\pi_{1}(B_{0}, a) \to \pi_{1}(B, a)\) 是拓扑群同构。

\begin{enumerate}
\def\labelenumi{\arabic{enumi})}
\setcounter{enumi}{1}
\item
设 B 是拓扑空间。设 \(b\) 和 \(b'\) 是属于同一道路连通分支的 B 的点，并设 \(\gamma\) 是 \(\varpi_{b,b'}(B)\) 的元素。由容许子群的定义，对于 \(\pi_1(B,b')\) 的每个容许子群 H，子群 \(\gamma H\gamma^{-1}\) 是 \(\pi_1(B,b)\) 的容许子群。因此，同构 \(u_\gamma \colon \delta \mapsto \gamma \delta \gamma^{-1}\) 从 \(\pi_1(B,b')\) 到 \(\pi_1(B,b)\)（参见 III，第 292 页）在这些群赋予容许拓扑时是同胚。
\item
  设 A 和 B 是拓扑空间，\(a\) 是 A 的一点，\(f\colon\mathrm{A}\to\mathrm{B}\) 是连续映射。置 \(b=f(a)\)。若 H 是 \(\pi_{1}(\mathrm{B},b)\) 的容许子群，则它在同态 \(\pi_{1}(f,a)\) 下的逆像是 \(\pi_{1}(\mathrm{A},a)\) 的容许子群。因此，当这些群赋予容许拓扑时，群同态 \(\pi_{1}(f,a)\colon\pi_{1}(\mathrm{A},a)\to\pi_{1}(\mathrm{B},b)\) 连续。
\end{enumerate}

命题 5. --- 设 B 是局部道路连通的拓扑空间，a 是 B 的一点。要使 \(\pi_1(\mathrm{B},a)\) 的子群 H 容许，必要且充分条件是存在 B 的覆叠 (E,p) 以及一点 \(x\in \mathrm{E}_a\)，使得 \(\mathrm{H} = p_{*}(\pi_{1}(\mathrm{E},x))\)。

设 \((\mathrm{E},p)\) 是 B 的覆叠，设 x 是 \(E_{a}\) 的一点；置 \(\mathrm{H}=p_{*}(\pi_{1}(\mathrm{E},x))\)，并证明它是 \(\pi_{1}(\mathrm{B},a)\) 的容许子群。设 b 是 B 的一点，V 是 b 的邻域，并且 \(\mathrm{E}_{\mathrm{V}}=(p^{-1}(\mathrm{V}),p_{\mathrm{V}})\) 是 V 的可平凡化覆叠。设 \(\gamma\in\varpi_{a,b}(\mathrm{B})\)；证明对于每个 \(\delta\in\pi_{1}(\mathrm{V},b)\)，道路类 \(\gamma\delta\gamma^{-1}\) 属于 H。

由 III，第 301 页命题 3，存在唯一的道路严格同伦类 \(\gamma'\)，其起点为 E 中的 \(x\)，且 \(p_*(\gamma') = \gamma\)。设 \(y\) 是 \(\gamma'\) 的终点；有 \(p(y) = b\)（III，第 302 页，命题 4）。再令 \(\delta'\) 为 E 中起点为 \(y\) 且满足 \(p_*(\delta') = \delta\) 的唯一道路类。由于 \(\mathrm{E_V}\) 可平凡化，\(\delta'\) 是以 \(y\) 为基点的圈的类。于是，\(\gamma'\delta'(\gamma')^{-1}\) 是 E 中以 \(a\) 为基点的圈的类，其在 \(p_*\) 下的像是类 \(\gamma\delta\gamma^{-1}\)，这正是要证明的。

反之，设 H 是 \(\pi_1(B, a)\) 的容许子群。令 \(\lambda_a(B)\) 为对道路严格同伦关系取商所得的空间 \(\Lambda_a(B)\) 的商集，其中道路的起点为 \(a\)；由于严格同伦的道路具有相同的终点，终点映射 \(e: \Lambda_a(B) \to B\) 经取商定义了映射 \(\varepsilon: \lambda_a(B) \to B\)。道路类的复合赋予集合 \(\lambda_a(B)\) 以群 \(\pi_1(B, a)\) 的左作用。将此作用限制为群 H 的作用，并将其轨道集合记为 \(H\backslash\lambda_a(B)\)。

映射 \(\varepsilon\colon\lambda_{a}(B)\to B\) 经取商诱导出映射 \(q\colon H\backslash\lambda_{a}(B)\to B\)。道路类的复合赋予集合 \(H\backslash\lambda_a(B)\) 以群胚 \(\varpi(B)\) 关于映射 \(q\) 的右作用。

此作用无局部单值性。事实上，设 \(b\) 是 B 的一点，V 是 \(b\) 的一个邻域，使得 \(\gamma i_{*}(\pi_{1}(V,b))\gamma^{-1}\subset H\) 对每个道路类 \(\gamma\) 都成立，其中这些道路连接 B 中的 \(a\) 与 \(b\)，而 \(i\) 表示 V 到 B 的嵌入。由于 B 局部道路连通，还可进一步假设 V 道路连通。设 \(c\in V\)，令 \(\delta\) 为 \(\pi_1(B,c)\) 中的、以 \(c\) 为基点且像包含于 V 的圈的类，并令 \(\delta'\) 为 \(\lambda_a(B)\) 中满足 \(\varepsilon (\delta ') = c\) 的元素。令 \(\delta''\) 为 \(\varpi_{c,b}(V)\) 的元素，并置 \(\gamma = \delta '\delta ''\)。根据 V 的定义，\(\delta '\delta (\delta ')^{-1} = \gamma ((\delta '')^{-1}\delta \delta '')\gamma^{-1}\) 是 \(\pi_1(B,a)\) 中的元素，并属于 H。于是，\(\mathrm{H}\delta^{\prime}\cdot \delta = \mathrm{H}\delta^{\prime}\)；这证明 \(\pi_1(V,c)\) 在集合 \(q^{-1}(c)\) 上的作用是平凡的。

由 III，第 313 页命题 3，\(\mathrm{H}\backslash\lambda_{a}(\mathrm{B})\) 上存在唯一拓扑，使得 \(q\) 连续，B-空间 \((\mathrm{H}\backslash\lambda_{a}(\mathrm{B}),q)\) 是覆叠，并且 \(\varpi(\mathrm{B})\) 在此覆叠上的典范作用就是上面定义的作用。根据 III 第 305 页定理 1 的断言 b)，群 \(q_{*}(\pi_{1}(\mathrm{H}\backslash\lambda_{a}(\mathrm{B}),\mathrm{H}))\) 等于 H。

称群 \(\pi_1(B, b)\) 在集合 X 上的作用为容许的，如果典范映射 \(\pi_1(B, b) \to \mathfrak{S}_{\mathrm{X}}\) 的核是 \(\pi_1(B, b)\) 的开子群。这等价于说，同态 \(\pi_1(B, b) \to \mathfrak{S}_{\mathrm{X}}\) 在群 \(\mathfrak{S}_{\mathrm{X}}\) 赋予离散拓扑时是连续的。于是，当 X 赋予离散拓扑时，映射 \(\pi_1(B, b) \times X \to X\) 连续。反之，考虑 \(\pi_{1}(B,b)\) 在离散空间 X 上的连续作用。设 x 是 X 的一点；根据假设，它的稳定子是 \(\pi_{1}(B,b)\) 的开子群 H。对于 \(\gamma\in\pi_{1}(B,b)\)，\(\gamma\cdot x\) 的固定子是子群 \(\gamma H\gamma^{-1}\)，因此 \(\pi_{1}(B,b)\) 中固定 x 的轨道的每个元素的子群，是各子群 \(\gamma H\gamma^{-1}\) 的交（其中 \(\gamma\) 遍历 \(\pi_{1}(B,b)\)）。它是开子群，因为 \(\pi_{1}(B,b)\) 的正规开子群构成其拓扑的基（III，第 315 页，命题 4）。这说明，\(\pi_{1}(B,b)\) 在离散空间 X 上的连续作用，只要是传递的，或者更一般地说，只要 X 中各元素的轨道集合是有限的，就是容许的。

对于任意离散群 \(G\)、任意以 \(G\) 为群的主覆叠 \(E\) 在 \(B\) 上，以及任意点 \(x \in E_{b}\)，回顾映射 \(h_{(\mathrm{E},x)}\) 从 \(\pi_{1}(\mathrm{B},b)\) 到 G，它把 \(\gamma \in \pi_{1}(\mathrm{B},b)\) 关联到 G 中的唯一元素 \(g \in G\)，该元素满足 \(x \cdot g = x \cdot \gamma^{-1}\)，是群同态（III，第 306 页，命题 5）。

命题 6. --- 设 B 是拓扑空间，b 是 B 的一点。给群 \(\pi_1(B, b)\) 赋予容许拓扑。

\begin{enumerate}
\def\labelenumi{\alph{enumi})}
\item
  对于任意离散群 G、以 G 为群的 B 的主覆叠 E 以及任意 \(x \in E_{b}\)，同态 \(h_{(E,x)}\) 是拓扑群的连续同态。
\item
  对于 B 的每个覆叠 E，\(\pi_{1}(B,b)\) 在纤维 \(E_{b}\) 上的典范作用都是容许的。
\end{enumerate}

只需证明第二个断言。令 K 为满足 \(k \in \pi_{1}(B, b)\) 且 \(x \cdot k = x\) 对所有 \(x \in E_{b}\) 都成立的元素 k 所成的集合。证明 K 是 \(\pi_{1}(B, b)\) 的容许子群。

设 \(b'\) 是 B 的一点，\(V'\) 是 B 中 \(b'\) 的邻域，在其上方覆叠 E 可平凡化。设 \(\gamma\) 是 B 中连接 \(b\) 与 \(b'\) 的道路 c 的类，\(c'\) 是 E 中从 x 出发且提升 c 的唯一道路。对于所有 \(\delta \in \pi_1(V', b')\) 和每个 \(x' \in E_{b'}\)，有 \(x' \cdot \delta = x'\)。因此，对于 \(x \in E_b\) 和 \(\delta \in \pi_1(V', b')\)，有

\[
x \cdot \gamma \delta \gamma^ {- 1} = ((x \cdot \gamma) \cdot \delta) \cdot \gamma^ {- 1} = (x \cdot \gamma) \cdot \gamma^ {- 1} = x,
\]

所以 \(\gamma\delta\gamma^{-1}\) 属于 K。换言之，K 是容许子群，命题得证。

命题 7. --- 设 B 是连通且局部道路连通的拓扑空间，b 是 B 的一点。给群 \(\pi_1(B, b)\) 赋予容许拓扑。

\begin{enumerate}
\def\labelenumi{\alph{enumi})}
\item
  设 G 是离散拓扑群。对于每个连续同态 \(f\colon \pi_{1}(\mathrm{B},b)\to \mathrm{G}\)，存在以 G 为群的 B 的主覆叠 E 以及一点 \(x\) 属于 \(\mathrm{E}_{b}\)，使得 \(h_{(\mathrm{E},x)} = f\)。
\item
  对于每个赋有群 \(\pi_{1}(B,b)\) 的容许右作用的离散拓扑空间 F，存在 B 的覆叠 E，使得 \(\pi_{1}(B,b)\)-集 F 和 \(E_{b}\) 同构。
\end{enumerate}

证明 \(a\)。设 \(f\) 是从 \(\pi_1(B, b)\) 到离散群 G 的连续同态。其核是 \(\pi_1(B, b)\) 的开正规子群 K；令 H 为群 \(\pi_1(B, b)/K\)，令 \(\overline{f} \colon H \to G\) 为 \(f\) 经取商诱导的群同态。由 III，第 316 页命题 5，存在 B 的连通覆叠 E'，使得子群 K 是纤维 \(E'_{b}\) 的每一点的稳定子；覆叠 E' 是以 H = \(\pi_1(B, b)/K\) 为群的主覆叠。设 \(x'\) 是 \(E'_{b}\) 的一点；同态 \(h_{(E', x')} \colon \pi_1(B, b) \to H\) 是满射且核为 K（III，第 306 页，命题 5）。再令 \(\varphi \colon H \to G\) 为唯一的群同态，使得 \(\varphi \circ h_{(E', x')} = f\)。B 的相伴覆叠 \(E = E' \times^{H} G\) 是以 G 为群的主覆叠，并且有 \(h_{(E, x)} = \varphi \circ h_{(E', x')} = f\)（III，第 307 页，例 2）。断言 \(a\) 得证。

证明 \(b\)。设 F 是赋有群 \(\pi_1(B,b)\) 的容许右作用的集合。记 \(f: \pi_1(B,b) \to \mathfrak{S}_{\mathrm{F}}\) 为此作用，并给群 \(\mathfrak{S}_{\mathrm{F}}\) 赋予离散拓扑。由 \(a\)，存在以 \(\mathfrak{S}_{\mathrm{F}}\) 为群的 B 的主覆叠 E 以及一点 \(x \in \mathrm{E}\)，使得同态 \(h_{(\mathrm{E},x)}\) 等于 \(f\)。群 \(\pi_1(B,b)\) 在点 \(b\) 上的相伴覆叠 \(\mathrm{E} \times^{\mathfrak{S}_{\mathrm{F}}} \mathrm{F}\) 的纤维，被识别为 \(\pi_1(B,b)\) 在 F 上的作用（III，第 306 页，例 1）。

注记。--- 4) \(\pi_1(B, b)\) 上的容许拓扑是使 \(\pi_1(B, b)\) 在每个 B 的连通覆叠 E 的 \(E_b\) 上的作用连续的最粗拓扑（III，第 318 页，命题 6 和第 318 页，命题 7）。若 E 是 B 的连通覆叠，且 \(x\) 是纤维 \(E_b\) 的一点，则映射 \(c' \mapsto c'(1)\) 从 \(\Lambda_x(E)\) 到 E 连续。因此，映射 \(c \mapsto x \cdot c\) 从 \(\Lambda_b(B)\) 到 E 连续（III，第 302 页，命题 3 的推论 2）。通过限制到 \(\Omega_b(B)\) 并取商，得到映射 \(\gamma \mapsto x \cdot \gamma\)，当给群 \(\pi_1(B, b)\) 赋予 \(\Omega_b(B)\) 的商拓扑时该映射连续。因此，\(\pi_{1}(B,b)\) 上的容许拓扑比紧收敛拓扑的商拓扑更粗。它可以严格更粗（III，第 337 页，习题 7）。

\begin{enumerate}
\def\labelenumi{\arabic{enumi})}
\setcounter{enumi}{4}
\tightlist
\item
设 B 是连通且局部道路连通的拓扑空间，\(a\) 是 B 的一点，H 是 \(\pi_1(B,a)\) 的子群。由上面的注记，若 H 容许，则它对于紧收敛拓扑的商拓扑也是开子群。反之，设 H 是 \(\pi_1(B,a)\) 的正规子群，且对于紧收敛拓扑的商拓扑为开子群。设 \(x\) 是 B 的一点，\(\gamma \in \varpi_{a,x}(B)\)。此时 \(\gamma^{-1}H\gamma\) 是 \(\pi_1(B,x)\) 的子群，对紧收敛拓扑的商拓扑仍是开子群（III，第 293 页，注记 3）。于是存在 B 中 \(x\) 的邻域 V，使得 \(\gamma^{-1}H\gamma\) 包含每个像包含于 V 的、以 \(x\) 为基点的圈的类。因此有 \(\gamma \pi_1(V,x)\gamma^{-1} \subset H\)；由于 H 正规，这说明 H 是 \(\pi_1(B,a)\) 的容许子群。
\end{enumerate}

